% Solutions to Exercises  (full solutions, all chapters)
% Add one \section per chapter as chapters are written. Reference each exercise
% with \cref{ex:...} so numbering stays automatic.
\chapter{Solutions to Exercises}
\label{app:D}

This appendix gives full solutions to every exercise. Solutions are grouped by
chapter and reference each exercise by its number.

\section{Chapter 1}
\label{sec:soln-ch1}

\paragraph{\Cref{ex:intro-thirdpoint}.}
Write the line and the cubic in homogeneous coordinates on \(\CP^2\).
Substituting the linear equation of the line into the cubic gives a single
homogeneous cubic equation in one projective parameter along the line, i.e.\ a
degree-three form. Over \(\C\), and counting multiplicity, such a form has
exactly three roots, so the line meets the cubic in three points counted with
multiplicity. Two of them are \(P\) and \(Q\); the third is \(R\). Passing to
the projective plane is what guarantees the count is exact --- it supplies the
points at infinity where, for example, a vertical line meets the curve. The
precise statement is B\'ezout's theorem (\cref{ch:lines_divisors}).

\paragraph{\Cref{ex:intro-tangent}.}
As \(Q \to P\) along the curve, the secant line through \(P\) and \(Q\)
approaches the tangent line at \(P\). The chord construction needs two
intersection points to start from; when the two points coincide, ``meeting the
curve at \(P\) twice'' is exactly the condition that the line be tangent there,
i.e.\ that it meet \(E\) at \(P\) with intersection multiplicity at least two.
So doubling \(P \oplus P\) uses the tangent at \(P\), whose third intersection
with \(E\) is then reflected. Intersection multiplicity is made precise in
\cref{ch:lines_divisors}, and the tangent (doubling) law in
\cref{ch:chord_tangent_cases}.

\paragraph{\Cref{ex:intro-liouville}.}
A holomorphic function on \(\C/\Lambda\) lifts to a holomorphic function on
\(\C\) that is periodic with respect to \(\Lambda\); its values are determined
by those on the (compact) fundamental parallelogram, so it is bounded, hence by
Liouville's theorem constant. Therefore the only holomorphic functions on
\(\C/\Lambda\) are the constants (\cref{ch:elliptic_liouville}).

\paragraph{\Cref{ex:intro-map-arrows}.}
Reading the arrows of \cref{fig:dependency}:
\begin{itemize}
\item \(T \xrightarrow{\ \Phi\ } (\wp,\wp')\): the Weierstrass functions are
  built on the torus in \cref{ch:weierstrass_wp}.
\item \((\wp,\wp') \to E\): the differential equation embeds the torus as the
  cubic, \cref{ch:weierstrass_wp,ch:embedding}.
\item \(E \to (\text{a line cuts } P+Q+R)\): intersecting \(E\) with a line,
  \cref{ch:lines_divisors}.
\item \((\text{line}) \to (\text{principal divisors have Abel sum zero})\):
  a line gives a principal divisor, \cref{ch:lines_divisors}.
\item \((\text{Abel sum zero}) \to T\) with label \(P+Q+R=O\): Abel's sum law,
  \cref{ch:abel}, transported in \cref{ch:transport}.
\end{itemize}

\section{Chapter 2}
\label{sec:soln-ch2}

\paragraph{\Cref{ex:lat-area}.}
Writing \(\omega_1 = x_1 + iy_1\) and \(\omega_2 = x_2 + iy_2\), the
parallelogram spanned by the vectors \((x_1,y_1)\) and \((x_2,y_2)\) has area
\(|x_1 y_2 - x_2 y_1|\). A short computation gives
\(\operatorname{Im}(\overline{\omega_1}\,\omega_2)
= x_1 y_2 - y_1 x_2\), so the area is
\(|\operatorname{Im}(\overline{\omega_1}\,\omega_2)|\), which is nonzero exactly
because \(\omega_1,\omega_2\) are \(\R\)-independent. Under a change of basis by
\(M \in \mathrm{GL}_2(\Z)\) (\cref{prop:lat-basis-change}) the area scales by
\(|\det M| = 1\), so it is a lattice invariant, the \emph{covolume} of
\(\Lambda\).

\paragraph{\Cref{ex:lat-discrete}.}
The function \((s,t) \mapsto |s\omega_1 + t\omega_2|\) is continuous and, by
\(\R\)-independence, strictly positive on the compact set
\(K = \{(s,t) \in \R^2 : \max(|s|,|t|) = 1\}\); let \(c > 0\) be its minimum.
For integers \(a,b\) not both zero, put \(n = \max(|a|,|b|) \ge 1\); then
\((a/n, b/n) \in K\), so \(|a\omega_1 + b\omega_2| = n\,|(a/n)\omega_1 +
(b/n)\omega_2| \ge cn = c\,\max(|a|,|b|)\). In particular
\(m(\Lambda) \ge c > 0\). A bounded region \(\{|\lambda| \le R\}\) then forces
\(\max(|a|,|b|) \le R/c\), leaving finitely many \((a,b)\), so finitely many
lattice points.

\paragraph{\Cref{ex:lat-square}.}
For \(\Lambda = \Z + \Z i\), the matrix
\(M = \left(\begin{smallmatrix} 0 & -1 \\ 1 & 0 \end{smallmatrix}\right)
\in \mathrm{SL}_2(\Z)\) sends \(\tau = i\) to
\(\frac{0\cdot i - 1}{1 \cdot i + 0} = -1/i = i\), so it fixes \(i\). It
corresponds to the basis change \((\omega_1,\omega_2) = (1,i) \mapsto
(\omega_1', \omega_2') = (i, -1) = (i\cdot 1,\, i\cdot i)\), i.e.\ multiplication
of the lattice by \(i\): the order-four rotational symmetry of the square
lattice. (Its square \(-I\) acts trivially on \(\tau\), reflecting that \(-1\)
also preserves \(\Lambda\).)

\paragraph{\Cref{ex:lat-homothety-iso}.}
If \(z \sim w\) in \(\C/\Lambda\), i.e.\ \(z - w \in \Lambda\), then
\(\alpha z - \alpha w \in \alpha\Lambda = \Lambda'\), so
\([z] \mapsto [\alpha z]\) is well defined; the inverse is
\([w] \mapsto [\alpha^{-1} w]\), so it is a bijection, and it is continuous both
ways because \(\pi' \circ (\alpha \cdot) = (\text{map}) \circ \pi\) with
\(z \mapsto \alpha z\) continuous. In the charts of
\cref{thm:lat-riemann-surface} the map lifts to \(z \mapsto \alpha z\), which is
holomorphic with holomorphic inverse; hence it is biholomorphic.

\paragraph{\Cref{ex:lat-chart-radius}.}
Take \(\Lambda = \Z + \Z i\), so \(m(\Lambda) = 1\) and \(m(\Lambda)/2 = 1/2\).
For any \(r > 1/2\) the open disk \(D(0, r)\) contains both \(-1/2\) and
\(1/2\), and these differ by \(1 \in \Lambda\); hence
\(\pi(-1/2) = \pi(1/2)\) while \(-1/2 \ne 1/2\), so \(\pi\) is not injective on
\(D(0,r)\). Conversely, for \(r \le 1/2\) any two points \(z, z' \in D(z_0, r)\)
satisfy \(|z - z'| < 2r \le 1 = m(\Lambda)\); a nonzero lattice vector has length
at least \(m(\Lambda)\), so \(z - z' \in \Lambda\) forces \(z = z'\) and \(\pi\)
is injective. Thus injectivity holds exactly up to radius \(m(\Lambda)/2\), and
the threshold is sharp. (The proof's choice \(r < m(\Lambda)/2\) is a convenient
strict bound; any \(r \le m(\Lambda)/2\) would serve.)

\section{Chapter 3}
\label{sec:soln-ch3}

\paragraph{\Cref{ex:tg-lie-holo}.}
Cover \(T \times T\) by products \(U_{z_0} \times U_{w_0}\) of the charts
\(\varphi_{z_0}, \varphi_{w_0}\) from \cref{thm:lat-riemann-surface}. In the
source product chart the points are \((z, w)\) with \(z \in D(z_0, r)\),
\(w \in D(w_0, r)\), and \(m([z],[w]) = [z+w]\); reading the target through the
chart \(\varphi_{z_0 + w_0}\) (valid near \([z_0+w_0]\)) the map is
\((z, w) \mapsto z + w\), a holomorphic function of two complex variables. Since
\(\pi\) is a local biholomorphism, \(m\) is holomorphic. For \(\iota\), the
local expression through \(\varphi_{z_0}\) and \(\varphi_{-z_0}\) is
\(z \mapsto -z\), holomorphic with holomorphic inverse \(z \mapsto -z\); hence
\(\iota\) is biholomorphic.

\paragraph{\Cref{ex:tg-two-torsion}.}
A point \([z]\) is fixed by \(\iota\) iff \([-z] = [z]\), i.e.\ \(2z \in
\Lambda\), i.e.\ \(z \in \tfrac12\Lambda\). Modulo \(\Lambda\), the quotient
\(\tfrac12\Lambda / \Lambda \cong (\tfrac12\Z/\Z)^2 \cong (\Z/2\Z)^2\) has
exactly four elements, represented by
\(0,\ \tfrac{\omega_1}{2},\ \tfrac{\omega_2}{2},\ \tfrac{\omega_1+\omega_2}{2}\).
These are precisely the solutions of \(2[z] = [0]\), the two-torsion subgroup
\(T[2] \cong (\Z/2\Z)^2\).

\paragraph{\Cref{ex:tg-first-iso}.}
The map \(e : \R \to S^1\), \(t \mapsto e^{2\pi i t}\), satisfies
\(e(t_1 + t_2) = e(t_1)e(t_2)\), so it is a homomorphism; it is continuous and
surjective (every unit-modulus number is \(e^{2\pi i t}\) for some \(t\)), and
\(e(t) = 1 \iff t \in \Z\), so \(\ker e = \Z\). The first isomorphism theorem
gives a continuous group isomorphism \(\bar e : \R/\Z \to S^1\). Finally
\(\R/\Z\) is the continuous image of the compact interval \([0,1]\), hence
compact, and \(S^1\) is Hausdorff, so the continuous bijection \(\bar e\) is a
homeomorphism.

\paragraph{\Cref{ex:tg-euler-refine}.}
Cut \(P_\Lambda\) along a diagonal into two triangles before gluing. After
identification one finds, for example, \(V = 1\) (all corners still collapse to
one vertex), \(E = 3\) (the two boundary loops \(a, b\) plus the diagonal
\(d\)), and \(F = 2\) (the two triangles). Then
\(\chi = V - E + F = 1 - 3 + 2 = 0\), the same value: the Euler characteristic
is independent of the chosen CW structure.

\paragraph{\Cref{ex:tg-homeo-not-biholo}.}
Both lattices are of the form \(\Z + \Z\tau\) with \(\tau = i\) and
\(\tau = 2i\), so by \cref{cor:tg-abelian} each torus is homeomorphic (and
diffeomorphic) to \(S^1 \times S^1\); composing the two homeomorphisms gives a
homeomorphism between them. For the complex structure, \cref{prop:lat-normal-form}
says \(\C/(\Z+\Z i)\) and \(\C/(\Z+\Z\,2i)\) are biholomorphic iff \(i\) and
\(2i\) are related by an element of \(\mathrm{SL}_2(\Z)\). Granting that both lie
in the closed fundamental domain and are not identified there, they are
inequivalent, so the tori are not biholomorphic. There is no contradiction with
\cref{cor:tg-abelian}: that corollary asserts only a \emph{topological}
isomorphism, and \cref{warn:tg-not-biholo} stresses that the homeomorphism
\(\Psi\) is not holomorphic. Homeomorphic surfaces can carry inequivalent complex
structures.

\section{Chapter 4}
\label{sec:soln-ch4}

\paragraph{\Cref{ex:ell-field}.}
For \(f, g\) elliptic, each of \(f \pm g, fg, f/g\) is meromorphic, and
periodicity is pointwise: \((fg)(z+\lambda) = f(z+\lambda)g(z+\lambda) =
f(z)g(z)\), and similarly for the others; the constants \(0, 1\) are elliptic, so
we have a field. Differentiating \(f(z+\lambda) = f(z)\) gives
\(f'(z+\lambda) = f'(z)\), so \(f'\) is elliptic. For fixed \(a\), \((f(\cdot+a))
(z+\lambda) = f(z+a+\lambda) = f(z+a)\), so \(f(z+a)\) is elliptic. If
\(f(-z) = f(z)\), differentiate: \(-f'(-z) = f'(z)\), i.e.\ \(f'(-z) = -f'(z)\),
so \(f'\) is odd.

\paragraph{\Cref{ex:ell-no-order-one}.}
An order-one function would have, in a period parallelogram, a single simple
pole and no other poles; the residue there is nonzero (the defining feature of a
simple pole), contradicting \(\sum \Res = 0\) (\cref{thm:ell-residues}). Hence no
order-one elliptic function exists. For order two the total pole multiplicity is
\(2\): either two simple poles at \(p_1 \ne p_2\) with residues \(r_1, r_2\), and
\cref{thm:ell-residues} forces \(r_1 + r_2 = 0\); or a single pole of order two,
with Laurent expansion \(c_{-2}(z-p)^{-2} + c_{-1}(z-p)^{-1} + \cdots\) whose
residue is \(c_{-1}\), and since it is the only pole, \(\sum\Res = c_{-1} = 0\).

\paragraph{\Cref{ex:ell-double-pole-zeros}.}
By \cref{thm:ell-valence} the order-two function \(f\) has exactly two zeros
\(z_1, z_2\) (with multiplicity). Its only pole is the double pole at \([0]\), so
\cref{thm:ell-abel} reads \(z_1 + z_2 - (0 + 0) \in \Lambda\), i.e.\
\(z_2 \equiv -z_1 =: -w\). The zeros are \([w], [-w]\). Now \(f - f(z_0)\) is
again elliptic of order two with the same double pole at \([0]\); one zero is
\(z_0\), so by the same argument its second zero is \(-z_0\). Thus \(f\) takes
each value at a pair \(\pm z_0\), which is the hallmark of an even function
\(f(-z) = f(z)\).

\paragraph{\Cref{ex:ell-sum-values}.}
The function \(f - c_1\) is elliptic of order \(N\) with the same poles
\(p_1, \dots, p_N\) as \(f\) and with zeros the solutions \(z_1^{(1)}, \dots\) of
\(f = c_1\). By \cref{thm:ell-abel}, \(\sum_k z_k^{(1)} - \sum_k p_k \in
\Lambda\). The same for \(c_2\): \(\sum_k z_k^{(2)} - \sum_k p_k \in \Lambda\).
Subtracting the two congruences, the pole sums cancel and
\(\sum_k z_k^{(1)} \equiv \sum_k z_k^{(2)} \pmod \Lambda\). Hence the sum of the
solution set of \(f = c\) is, modulo \(\Lambda\), independent of \(c\); its common
value (equal to \(\sum p_k\) mod \(\Lambda\)) is an invariant of \(f\).

\paragraph{\Cref{ex:ell-winding}.}
Let \(\gamma(s) = a + s\omega_1\), \(s \in [0,1]\), avoid the zeros and poles of
\(f\). Then \(\int_a^{a+\omega_1} (f'/f)\,dz = \int_0^1 \frac{d}{ds}\log f(\gamma(s))\,ds
= \log f(a+\omega_1) - \log f(a)\) for a continuous branch of \(\log f\) along
\(\gamma\). The real part is \(\log|f(a+\omega_1)| - \log|f(a)| = 0\) because
\(f(a+\omega_1) = f(a)\); the imaginary part is the net change of argument,
\(2\pi k_1\), where \(k_1 \in \Z\) is the number of times the image curve
\(s \mapsto f(\gamma(s))\) winds around \(0\). Hence the integral equals
\(2\pi i\,k_1\), and \(k_1\) is precisely that winding number.

\section{Chapter 5}
\label{sec:soln-ch5}

\paragraph{\Cref{ex:div-hom}.}
By \cref{prop:div-prin-subgroup}, \(\divisor(fg) = \divisor(f) + \divisor(g)\),
so \(\divisor\) is a homomorphism from \((\mathcal{M}(X)^\times, \cdot)\) to
\((\Div(X), +)\). If \(\divisor(f) = 0\) then \(f\) has no zeros or poles, so it
is holomorphic and nonvanishing on the compact \(X\); by \cref{thm:ell-const}
(or the maximum modulus principle) \(f\) is a nonzero constant. Conversely
\(\divisor(c) = 0\). Hence \(\ker\divisor = \C^\times\).

\paragraph{\Cref{ex:div-ord-welldef}.}
Independence of coordinate is \cref{rem:div-ord-welldef}: a coordinate change
\(t = s\,h(s)\), \(h(0)\ne0\), turns \(f = t^m u\) into
\(f = s^m\,(h^m\, u\!\circ\! t)\) with a nonvanishing unit, so the exponent
\(m\) is unchanged. For the sum, write \(f = t^{a}u\), \(g = t^{b}v\) with
\(a = \ord_p f \le b = \ord_p g\) and units \(u, v\). Then
\(f + g = t^{a}(u + t^{b-a}v)\); the factor \(u + t^{b-a}v\) is holomorphic, so
\(\ord_p(f+g) \ge a = \min(a,b)\), and if \(a < b\) it equals \(u(0) \ne 0\) at
\(t=0\), giving \(\ord_p(f+g) = a\).

\paragraph{\Cref{ex:div-pic-split}.}
The map \(D \mapsto (\degdiv D,\ D - (\degdiv D)[0])\) is a homomorphism, since
\(\degdiv\) is and the second slot is linear in \(D\); its second component has
degree \(\degdiv D - \degdiv D = 0\), so it lands in \(\Z \oplus \Div^0(T)\). The
inverse is \((n, E) \mapsto E + n[0]\), so it is an isomorphism. Principal
divisors have degree \(0\) and lie in the \(\Div^0\) slot, so quotienting gives
\(\Pic(T) \cong \Z \oplus \Pic^0(T)\); with \(\Pic^0(T) \cong T\) from
\cref{thm:div-pic-iso}, \(\Pic(T) \cong \Z \oplus T\).

\paragraph{\Cref{ex:div-distinct-points}.}
If \([a] = [b]\) then \([a] - [b] = 0 = \divisor(1)\) is principal. Conversely, if
\([a] - [b] = \divisor(f)\) is principal, then it lies in \(\Div^0(T)\) and
\cref{thm:ell-abel} (via \cref{prop:div-sigma-kills-prin}) gives
\(\smap([a]-[b]) = [a] - [b] = 0\) in \(T\), i.e.\ \([a] = [b]\). On \(\CP^1\)
the function \((z-a)/(z-b)\) has divisor \([a] - [b]\) for any \(a \ne b\), so the
statement fails there. The torus version recovers \cref{cor:ell-order2}: a
principal divisor \([a] - [b]\) with \(a \ne b\) would be the divisor of an
order-one elliptic function (one simple zero, one simple pole), which does not
exist.

\paragraph{\Cref{ex:div-sigma-compute}.}
The divisor \(D = [\tfrac12] + [\tfrac{\tau}{2}] + [\tfrac{1+\tau}{2}] - 3[0]\)
has \(\degdiv D = 3 - 3 = 0\). Its Abel sum is
\[
  \smap(D) = \Bigl[\tfrac12 + \tfrac{\tau}{2} + \tfrac{1+\tau}{2}\Bigr]
  = \Bigl[\tfrac{2 + 2\tau}{2}\Bigr] = [\,1 + \tau\,] = [0],
\]
since \(1 + \tau \in \Lambda\). So \(D\) has degree zero and Abel sum zero, hence
is a candidate for a principal divisor --- and indeed it is: the three nonzero
half-periods \(\tfrac12, \tfrac{\tau}{2}, \tfrac{1+\tau}{2}\) are exactly the
two-torsion points of \cref{ex:tg-two-torsion}, and \(D = \divisor(\wp')\), the
divisor of the Weierstrass derivative (simple zeros at the half-periods, a
triple pole at \([0]\)), as we will see in \cref{ch:weierstrass_wp}.

\section{Chapter 6}
\label{sec:soln-ch6}

\paragraph{\Cref{ex:rr-point-lemma}.}
Let \(D = \sum n_q\,q\). An \(f \in L(D+p)\) may have a pole at \(p\) of order up
to \(n_p + 1\). In a coordinate \(t\) at \(p\), the map
\(\lambda : L(D+p) \to \C\) sending \(f\) to its Laurent coefficient of order
\(-(n_p+1)\) (the most negative now allowed) is linear. Its kernel is exactly the
functions whose pole order at \(p\) is at most \(n_p\), i.e.\ \(L(D)\). Hence
\(L(D+p)/L(D) \hookrightarrow \C\), so \(\ell(D+p) - \ell(D) = \dim\operatorname{im}\lambda
\in \{0,1\}\).

\paragraph{\Cref{ex:rr-neg-degree}.}
If \(\deg D < 0\) and \(0 \ne f \in L(D)\), then \(\divisor(f) \ge -D\) has
degree \(0\) (\cref{thm:div-degzero}) but \(\deg(-D) > 0\), impossible; so
\(L(D) = 0\). For the inequality, induct on \(\deg D\) downward using the point
lemma: any \(D\) can be written \(D = D_0 + \sum_{j=1}^{k} p_j\) with
\(\deg D_0 \le 0\) and \(k = \deg D - \deg D_0\); since
\(\ell(D_0) \le 1\) (it is \(0\) for negative degree, and for degree \(0\) a
nonzero \(f \in L(D_0)\) has \(\divisor(f) = -D_0\), so \(D_0 \sim 0\) and
\(\ell(D_0) = 1\)) and each added point raises \(\ell\) by at most
one (\cref{ex:rr-point-lemma}), \(\ell(D) \le 1 + (\deg D - \deg D_0) \le \deg D + 1\).

\paragraph{\Cref{ex:rr-degree-one-unique}.}
A divisor of degree one is \(\sum n_q\,q\) with \(\sum n_q = 1\); it is effective
(\(\ge 0\)) iff all \(n_q \ge 0\), which with sum \(1\) means exactly one
coefficient equals \(1\) and the rest \(0\), i.e.\ a single point. By
\cref{prop:rr-counts}, \(\ell([a]) = 1\), so \(L([a]) = \C\) (constants); the
effective divisors linearly equivalent to \([a]\) are the \(\divisor(f) + [a]\)
for \(0 \ne f \in L([a])\), i.e.\ \(f\) constant, giving only \([a]\) itself.

\paragraph{\Cref{ex:rr-canonical-deg}.}
By \cref{prop:rr-canonical}, \(K \sim 0\), so \(\ell(K) = \ell(0) = 1 = g\). Taking
\(D = K\) in Riemann--Roch: \(\ell(K) - \ell(K - K) = \deg K + 1 - g\), i.e.\
\(\ell(K) - \ell(0) = \deg K + 1 - g\). With \(\ell(K) = \ell(0) = 1\) and
\(g = 1\) this reads \(0 = \deg K + 1 - 1\), so \(\deg K = 0 = 2g - 2\),
consistent.

\paragraph{\Cref{ex:rr-cubic-count}.}
The pole orders at \(O\) are
\[
  1 \!:\! 0,\quad \wp \!:\! 2,\quad \wp^2 \!:\! 4,\quad \wp^3 \!:\! 6,\quad
  \wp' \!:\! 3,\quad \wp\wp' \!:\! 5,\quad \wp'^2 \!:\! 6,
\]
all \(\le 6\) and with no other poles, so all seven lie in \(L(6O)\). Since
\(\ell(6O) = 6 < 7\), they are linearly dependent: \(\sum c_i\,F_i = 0\)
nontrivially. The only two functions with a pole of order exactly \(6\) are
\(\wp^3\) and \(\wp'^2\); their leading Laurent coefficients at \(O\) are the
only order-\(6\) contributions, so for the order-\(6\) pole to cancel, the
coefficients of \(\wp^3\) and \(\wp'^2\) must be nonzero and balance. If either
were zero, the other would contribute an uncancelled order-six pole, contradicting
\(\sum c_i F_i = 0\) (a holomorphic, indeed zero, function). Hence both appear,
and solving for the relation gives \(\wp'^2 = 4\wp^3 - g_2\wp - g_3\).

\section{Chapter 7}
\label{sec:soln-ch7}

\paragraph{\Cref{ex:wp-convergence}.}
Lattice points are uniformly separated (\cref{ex:lat-discrete}), so disjoint disks
of radius \(\rho = \tfrac12 m(\Lambda)\) about them are contained in the disk of
radius \(R + \rho\) once \(|\lambda| \le R\); comparing areas bounds their number
by \((R+\rho)^2/\rho^2 = O(R^2)\). Splitting into annuli \(n \le |\lambda| < n+1\)
(each with \(O(n)\) points) gives \(\sum_\lambda |\lambda|^{-s} \le \sum_n O(n)\,
n^{-s} = \sum_n O(n^{1-s}) < \infty\) for \(s > 2\). The differentiated terms
\(\frac{d}{dz}\bigl((z-\lambda)^{-2} - \lambda^{-2}\bigr) = -2(z-\lambda)^{-3}\)
satisfy \(|{-2}(z-\lambda)^{-3}| \le 16/|\lambda|^3\) for \(|z| \le R\),
\(|\lambda| \ge 2R\), which is summable (\(s = 3\)); uniform convergence of the
differentiated series justifies term-by-term differentiation, and adding the
\(\lambda = 0\) term \(\tfrac{d}{dz}z^{-2} = -2z^{-3}\) gives
\(\wp'(z) = -2\sum_{\lambda \in \Lambda}(z-\lambda)^{-3}\).

\paragraph{\Cref{ex:wp-g2g3}.}
From \cref{prop:wp-laurent}, \(\wp = z^{-2} + 3G_4 z^2 + 5G_6 z^4 + O(z^6)\), so
\(a_2 = 3G_4\), \(a_4 = 5G_6\). Using \(\wp' = -2z^{-3} + 2a_2 z + 4a_4 z^3 +
\cdots\), one computes \(\wp'^2 - 4\wp^3 = -20a_2 z^{-2} - 28 a_4 + O(z^2)\).
Choosing \(g_2 = 20 a_2\) and \(g_3 = 28 a_4\) cancels the \(z^{-2}\) term and the
constant; substituting \(a_2 = 3G_4\), \(a_4 = 5G_6\) gives \(g_2 = 60G_4\),
\(g_3 = 140 G_6\).

\paragraph{\Cref{ex:wp-roots-distinct}.}
The cubic \(4x^3 - g_2 x - g_3 = 4\bigl(x^3 - \tfrac{g_2}{4}x - \tfrac{g_3}{4}\bigr)\)
is monic after dividing by \(4\), with \(p = -g_2/4\), \(q = -g_3/4\). The
discriminant of a monic cubic \(x^3 + px + q\) is \(\prod_{i<j}(e_i-e_j)^2 =
-4p^3 - 27q^2\). Hence
\[
  \prod_{i<j}(e_i-e_j)^2 = -4\Bigl(-\tfrac{g_2}{4}\Bigr)^3 - 27\Bigl(-\tfrac{g_3}{4}\Bigr)^2
  = \frac{g_2^3}{16} - \frac{27 g_3^2}{16} = \frac{g_2^3 - 27g_3^2}{16},
\]
so \(\Delta = g_2^3 - 27g_3^2 = 16\prod_{i<j}(e_i-e_j)^2\). By
\cref{prop:wp-roots}(3) the \(e_i\) are distinct, so \(\Delta \ne 0\).

\paragraph{\Cref{ex:wp-field}.}
Write \(f = f_+ + f_-\) with \(f_\pm(z) = \tfrac12(f(z) \pm f(-z))\), both
elliptic, \(f_+\) even and \(f_-\) odd. By the generation step of
\cref{thm:wp-generation}, even elliptic functions lie in \(\C(\wp)\), so
\(f_+ = A(\wp)\) for a rational \(A\). Since \(\wp'\) is odd, \(f_-/\wp'\) is even,
hence \(f_-/\wp' = B(\wp)\) for a rational \(B\); thus \(f_- = B(\wp)\,\wp'\) and
\(f = A(\wp) + B(\wp)\,\wp'\). The fact \(\wp' \notin \C(\wp)\) guarantees this
representation is genuine and unique: if
\(A_1(\wp) + B_1(\wp)\wp' = A_2(\wp) + B_2(\wp)\wp'\) then
\((B_1 - B_2)(\wp)\,\wp' = (A_2 - A_1)(\wp) \in \C(\wp)\) forces
\(B_1 = B_2\) (else \(\wp' \in \C(\wp)\)) and then \(A_1 = A_2\).

\paragraph{\Cref{ex:wp-recurrence}.}
Differentiating \(\wp'^2 = 4\wp^3 - g_2\wp - g_3\) gives \(2\wp'\wp'' =
(12\wp^2 - g_2)\wp'\), and dividing by \(2\wp'\) (not identically zero),
\(\wp'' = 6\wp^2 - \tfrac12 g_2\). Write \(\wp = z^{-2} + \sum_{n\ge1} c_n z^{2n}\)
with \(c_n = (2n+1)G_{2n+2}\). Then \(\wp'' = 6z^{-4} + \sum_{n\ge1}
2n(2n-1)c_n z^{2n-2}\), while
\(6\wp^2 = 6z^{-4} + 12\sum_{n\ge1} c_n z^{2n-2} + 6\sum_{n\ge2}\bigl(\sum_{i+j=n}
c_i c_j\bigr)z^{2n-2}\). Matching the coefficient of \(z^{2n-2}\) for \(n \ge 3\)
(the cases \(n=1,2\) absorb \(g_2\)) gives \(2n(2n-1)c_n = 12c_n +
6\sum_{i+j=n}c_ic_j\), i.e.
\[
  c_n = \frac{3}{(2n+3)(n-2)}\sum_{\substack{i+j=n\\ i,j\ge1}} c_i c_j
  \qquad (n \ge 3).
\]
Since \(c_1 = 3G_4\) and \(c_2 = 5G_6\), this recurrence determines every
\(c_n\), hence every \(G_{2n+2}\), as a polynomial in \(G_4\) and \(G_6\); for
instance \(c_3 = \tfrac{3}{9\cdot1}c_1^2 = \tfrac13 c_1^2\), giving
\(G_8 = \tfrac{3}{7}G_4^2\).

\paragraph{\Cref{ex:wp-sigma-quasi}.}
Because \(\zeta' = -\wp\) is periodic, \(g(z) := \zeta(z+\omega_k) - \zeta(z)\) has
\(g' = 0\), so \(g \equiv \eta_k\) is constant; at \(z = -\omega_k/2\), oddness of
\(\zeta\) gives \(\eta_k = \zeta(\omega_k/2) - \zeta(-\omega_k/2) = 2\zeta(\omega_k/2)\).
Since \((\log\sigma)' = \sigma'/\sigma = \zeta\), the function
\(\log\frac{\sigma(z+\omega_k)}{\sigma(z)}\) has derivative
\(\zeta(z+\omega_k) - \zeta(z) = \eta_k\), so it equals \(\eta_k z + c_k\) for some
constant \(c_k\), i.e.\ \(\sigma(z+\omega_k) = \sigma(z)\,e^{\eta_k z + c_k}\).
Evaluate at \(z = -\omega_k/2\): the left side is \(\sigma(\omega_k/2)\), the right
\(\sigma(-\omega_k/2)e^{-\eta_k\omega_k/2 + c_k} = -\sigma(\omega_k/2)
e^{-\eta_k\omega_k/2 + c_k}\) (oddness), and \(\sigma(\omega_k/2) \neq 0\) since
\(\omega_k/2 \notin \Lambda\); hence \(e^{c_k - \eta_k\omega_k/2} = -1\) and
\(\sigma(z+\omega_k) = -\sigma(z)\,e^{\eta_k(z+\omega_k/2)}\).

\paragraph{\Cref{ex:wp-converse-check}.}
Using \(\sigma(z+\omega_k - c) = -\sigma(z-c)\,e^{\eta_k(z - c + \omega_k/2)}\) for
each shift \(c \in \{a_i, b_i\}\),
\[
  \frac{F(z+\omega_k)}{F(z)}
  = \frac{\prod_{i=1}^n\bigl(-e^{\eta_k(z - a_i + \omega_k/2)}\bigr)}
         {\prod_{i=1}^n\bigl(-e^{\eta_k(z - b_i + \omega_k/2)}\bigr)} .
\]
The numerator is \((-1)^n e^{\eta_k(nz - \sum a_i + n\omega_k/2)}\) and the
denominator \((-1)^n e^{\eta_k(nz - \sum b_i + n\omega_k/2)}\); the factor
\((-1)^n\) cancels because the two products have the \emph{same number} of factors
(degree zero), and the \(nz\) and \(n\omega_k/2\) terms cancel, leaving
\(\exp\!\bigl(\eta_k(\sum b_i - \sum a_i)\bigr)\). This is \(1\) because
\(\sum a_i = \sum b_i\) (Abel sum zero). Thus degree zero is used to cancel
\((-1)^n\) and the \(z\)-dependence, and Abel sum zero to kill the residual
exponent. For \([a] - [0]\) with \(a \notin \Lambda\) the counts match (degree
zero), so the candidate \(\sigma(z-a)/\sigma(z-0)\) has the right zeros and poles,
but \(\sum a_i - \sum b_i = a \notin \Lambda\), so \(F(z+\omega_k)/F(z) =
e^{\eta_k a} \neq 1\) in general: \(F\) is not elliptic, consistent with there
being no order-one elliptic function (\cref{cor:ell-order2}).

\section{Chapter 8}
\label{sec:soln-ch8}

\paragraph{\Cref{ex:emb-holo-origin}.}
Near \([0]\), the homogeneous representative \((z^3\wp, z^3\wp', z^3)\) has middle
entry tending to \(-2 \ne 0\), so divide by it: in the chart \(Y = 1\) the map is
\(z \mapsto (X, Z) = (\wp/\wp',\, 1/\wp')\). From \(\wp = z^{-2} + O(z^2)\),
\(\wp' = -2z^{-3} + O(z)\): \(\wp/\wp' = \frac{z^{-2} + \cdots}{-2z^{-3} + \cdots}
= -\tfrac12 z + O(z^3)\) and \(1/\wp' = -\tfrac12 z^3 + O(z^5)\). Both are
holomorphic at \(z = 0\) with value \((0,0) = O\), and \(X = \wp/\wp'\) has
derivative \(-\tfrac12 \ne 0\) there, so it is a local holomorphic coordinate and
\(\Phi\) is an immersion at \(O\) (indeed \(z\) and \(X\) agree to first order up
to the factor \(-\tfrac12\)).

\paragraph{\Cref{ex:emb-flex}.}
Write the curve as \(G = Y^2 Z - 4X^3 + g_2 X Z^2 + g_3 Z^3 = 0\). At
\(O = [0:1:0]\), \(\nabla G = (-12X^2 + g_2 Z^2,\ 2YZ,\ Y^2 + 2g_2 XZ + 3g_3 Z^2)
= (0, 0, 1) \ne 0\), so \(O\) is smooth with tangent \(Z = 0\). Substituting
\(Z = 0\) into \(G\) gives \(-4X^3 = 0\), a triple root, so \(Z = 0\) meets \(E\)
at \(O\) with multiplicity three; \(O\) is a flex. Hence the line \(Z = 0\) cuts
the divisor \(3O\) on \(E\), so \(3O = \divisor(\text{line}/\text{line})\) is
principal --- concretely \(3O \sim 3O\), and any two lines cut linearly equivalent
divisors, making \(3O\) the class that will be the group identity.

\paragraph{\Cref{ex:emb-fiber-lemma}.}
\Cref{lem:emb-fibers} is proved in the text: \(\wp - \wp(w)\) has order two, with
zeros \(w, -w\), so \(\wp(z) = \wp(w) \Rightarrow z \equiv \pm w\). Thus the fibers
of \(\wp : T \to \CP^1\) are the orbits \(\{z, -z\}\) of \(z \mapsto -z\), giving a
bijection \(T/\{\pm1\} \to \CP^1\). The fixed points of \(z \mapsto -z\) are the
four two-torsion points (\cref{ex:tg-two-torsion}); there the two sheets come
together, so they map to the branch values, namely \(\wp(\omega_i/2) = e_i\) and
\(\wp(0) = \infty\).

\paragraph{\Cref{ex:emb-two-torsion-image}.}
At \(z = 0\), \(\Phi(0) = O\). At a half-period \(\omega_i/2\), \(\wp'(\omega_i/2)
= 0\) (\cref{prop:wp-roots}) and \(\wp(\omega_i/2) = e_i\), so
\(\Phi(\omega_i/2) = (e_i, 0)\). The points of \(E\) fixed by \((x,y) \mapsto
(x,-y)\) satisfy \(y = -y\), i.e.\ \(y = 0\) (the three \((e_i, 0)\)) or the point
at infinity \(O\); these are exactly the four images, matching the two-torsion.

\paragraph{\Cref{ex:emb-involution}.}
The map \(\iota_E(x,y) = (x,-y)\) preserves \(y^2 = 4x^3 - g_2 x - g_3\) (the
right side depends only on \(x\)), fixes \(O\), and satisfies \(\iota_E^2 = \id\),
so it is a holomorphic involution. Its invariant function is \(x = \wp\), and two
points share an \(x\)-value iff they are \(\pm\)-related, so the quotient
\(E/\iota_E\) is the \(x\)-line \(\CP^1\). Under \(\Phi\), \(\iota_E\) pulls back to
the map with \(\wp(z) \mapsto \wp(z)\), \(\wp'(z) \mapsto -\wp'(z) = \wp'(-z)\),
i.e.\ \(z \mapsto -z\). Since the double cover \(\wp : T \to \CP^1\) has deck group
\(\{\pm1\}\) of order two, its unique nontrivial deck transformation is
\(z \mapsto -z\), which is \(\iota_E\).

\section{Chapter 9}
\label{sec:soln-ch9}

\paragraph{\Cref{ex:abel-independence}.}
Take \(D_1 = [a] - [0]\) with \(a \notin \Lambda\): \(\deg D_1 = 0\) but
\(\smap(D_1) = [a] \ne 0\), so \(D_1\) is not principal (a principal divisor of
this shape would be an order-one elliptic function, excluded by
\cref{cor:ell-order2}). Take \(D_2 = 2[0]\): \(\smap(D_2) = 2[0] = 0\) but
\(\deg D_2 = 2 \ne 0\), so \(D_2\) is not principal. Thus each condition can hold
while the other fails, so neither alone characterizes principality.

\paragraph{\Cref{ex:abel-three-point}.}
\Cref{cor:abel-three-point} is immediate from \cref{thm:abel-law}:
\([p]+[q]+[r]-3[0]\) has degree \(0\), and its Abel sum is \([p]+[q]+[r]\), zero
iff \(p+q+r \in \Lambda\). For existence/uniqueness of \(R\): given \(P = [p]\),
\(Q = [q]\), the equation \(P + Q + R = O\) in the group \(T\) has the unique
solution \(R = -(P+Q) = [-(p+q)]\). Through the isomorphism \(\Phi\) this \(R\) is
a well-defined point of the cubic, the reflection of the third intersection.

\paragraph{\Cref{ex:abel-period-integral}.}
Label the sides of \(P_a\) as in \cref{sec:ell-contour}. Pairing side
\(\mathrm{I}\) with \(\mathrm{III}\) (translates by \(\omega_2\), reversed) and
using \(\omega_2\)-periodicity of \(F'/F\),
\(\int_{\mathrm{I}} + \int_{\mathrm{III}} = -\omega_2 \int_{a}^{a+\omega_1}
(F'/F)\,dz\); pairing \(\mathrm{II}\) with \(\mathrm{IV}\) gives
\(+\omega_1 \int_{a}^{a+\omega_2}(F'/F)\,dz\). Each remaining integral is the
change of a branch of \(\log F\) along a side; since \(F(a+\omega_k) = F(a)\), the
real part returns to its start and the imaginary part changes by \(2\pi\) times an
integer (the winding number, \cref{ex:ell-winding}), so each integral lies in
\(2\pi i\,\Z\). Hence \(\oint z(F'/F)\,dz = \omega_1\cdot 2\pi i\,k_2 - \omega_2
\cdot 2\pi i\,k_1 \in 2\pi i\,\Lambda\).

\paragraph{\Cref{ex:abel-addition-function}.}
The divisor \(D = [a] + [b] - [a+b] - [0]\) has \(\deg D = 2 - 2 = 0\) and Abel
sum \(\smap(D) = [a] + [b] - [a+b] - [0] = (a + b) - (a+b) = 0\) in \(T\). By
sufficiency in \cref{thm:abel-law}, \(D\) is principal, so some elliptic \(F\) has
\(\divisor(F) = D\); its zeros \(a, b\) and poles \(a+b, 0\) encode the
group sum \(a + b\). By contrast \([a] - [0]\) with \(a \ne 0\) has Abel sum
\([a] \ne 0\), so it is \emph{not} principal --- there is no elliptic function with
a single simple zero and a single simple pole, consistent with
\cref{cor:ell-order2}.

\paragraph{\Cref{ex:abel-kernel}.}
Both \(\deg\) and \(\smap\) are homomorphisms \(\Div(T) \to \Z\) and
\(\Div(T) \to T\), so \((\deg, \smap)\) is a homomorphism to \(\Z \oplus T\). It
is surjective: \((n, [a]) = (\deg, \smap)\bigl(n[0] + ([a] - [0])\bigr)\), since
that divisor has degree \(n\) and Abel sum \([a]\). Its kernel is the divisors
with \(\deg = 0\) and \(\smap = 0\), which by \cref{thm:abel-law} are exactly
\(\Prin(T)\). The first isomorphism theorem then gives \(\Div(T)/\Prin(T) =
\Pic(T) \cong \Z \oplus T\), matching \cref{ex:div-pic-split}.

\section{Chapter 10}
\label{sec:soln-ch10}

\paragraph{\Cref{ex:lin-vertical}.}
The vertical line \(x = c\) is \(X - cZ = 0\). On \(E\), \(y^2 = 4c^3 - g_2 c - g_3
=: y_0^2\), giving the two affine points \((c, \pm y_0)\); and the line passes
through \(O = [0:1:0]\) (it contains the direction \([0:1:0]\), since
\(X - cZ\) vanishes there). So \(L \cdot E = (c, y_0) + (c, -y_0) + O\), which has
degree three as required. By \cref{thm:lin-collinear} these three sum to \(O\):
\((c, y_0) \oplus (c, -y_0) \oplus O = O\), i.e.\ \((c, y_0) \oplus (c, -y_0) = O\).
Thus \((c, -y_0) = -(c, y_0)\), the reflection across the \(x\)-axis of
\cref{sec:emb-reflection}.

\paragraph{\Cref{ex:lin-mult-welldef}.}
If \(m, m'\) both satisfy \(m(P), m'(P) \ne 0\), then \(m/m'\) is a rational
function on \(E\) with neither zero nor pole at \(P\), hence
\(\ord_P(m/m') = 0\). Therefore \(\ord_P(\ell/m) = \ord_P(\ell/m') +
\ord_P(m'/m) = \ord_P(\ell/m')\), so the multiplicity does not depend on the
auxiliary form.

\paragraph{\Cref{ex:lin-tangent}.}
By definition \(I_P(L,E) = \ord_P(\ell/m)\ge 1\) exactly when \(P \in L\). Choose a
local coordinate \(t\) on \(E\) at \(P\) and parametrize \(E\) by \(\gamma(t)\) with
\(\gamma(0) = P\). Then \(I_P(L,E) \ge 2\) iff \((\ell/m)(\gamma(t))\) vanishes to
order \(\ge 2\), iff \(\frac{d}{dt}(\ell/m)(\gamma(t))|_0 = 0\), iff the line
\(\ell = 0\) contains the tangent direction \(\gamma'(0)\) --- i.e.\ \(L\) is the
tangent line to \(E\) at \(P\). When \(L\) is tangent, \(I_P(L,E) \ge 2\) absorbs
two of the three intersections (\cref{prop:lin-bezout}), leaving one further point
\(R'\) (possibly \(= P\) at a flex), so \(L \cdot E = 2P + R'\).

\paragraph{\Cref{ex:lin-infinity-mult}.}
With \(v = 4u^3 - g_2 u v^2 - g_3 v^3\) and \(v = O(u^3)\) near \(0\): substituting
the ansatz \(v = c_3 u^3 + c_5 u^5 + \cdots\) (only odd powers by the symmetry
\(u \mapsto -u\), \(v \mapsto -v\)), the terms \(g_2 u v^2\) and \(g_3 v^3\) are
\(O(u^7)\), so to leading order \(v = 4u^3 + O(u^5)\); hence \(c_3 = 4\) and
\(\ord_O(Z/Y) = \ord_O v = 3\). Since \(Y\) is a unit at \(O\), \(\ord_O(Z|_E) =
3\) and \(\divisor_0(Z|_E) = 3O\).

\paragraph{\Cref{ex:lin-three-collinear}.}
Let \(R\) be the third intersection of the line through \(P\) and \(Q\) with \(E\),
so \(L \cdot E = P + Q + R\). By \cref{thm:lin-collinear},
\(P \oplus Q \oplus R = O\), hence \(R = -(P \oplus Q)\) and
\(P \oplus Q = -R\), the reflection of \(R\). When \(P = Q\) the chord degenerates;
the correct line is the \emph{tangent} to \(E\) at \(P\), which meets \(E\) at
\(2P + R\) (\cref{ex:lin-tangent}), giving \(2P = -R\) --- the doubling rule.

\section{Chapter 11}
\label{sec:soln-ch11}

\paragraph{\Cref{ex:tr-unique-transport}.}
If \(\boxdot\) makes \(\Phi\) a homomorphism, then \(\Phi(u+v) = \Phi(u) \boxdot
\Phi(v)\); since \(\Phi\) is a bijection, every \(P = \Phi(u)\), \(Q = \Phi(v)\),
and \(P \boxdot Q = \Phi(u+v)\) is determined --- this is \(\oplus\). It is abelian
because \(u + v = v + u\) in \(\C\); \(P \oplus O = \Phi(u + 0) = \Phi(u) = P\);
and \(\Phi(u) \oplus \Phi(-u) = \Phi(0) = O\), so \(\ominus\Phi(u) = \Phi(-u)\).

\paragraph{\Cref{ex:tr-identity-case}.}
The line \(L = PO\) through \(P = (x,y)\) and \(O = [0:1:0]\) is the vertical line
\(x = x_P\) (the point at infinity \(O\) lies on every vertical line). It meets
\(E\) at \((x, y)\), \((x, -y) = \ominus P\), and \(O\); so \(L \cdot E = P +
\ominus P + O\). By \cref{thm:lin-collinear}, \(P \oplus \ominus P \oplus O = O\),
which is automatic, and reading the chord--tangent recipe: the third point of
\(PO\) is \(\ominus P\), and \(P \oplus O = \ominus(\ominus P) = P\). Hence \(O\) is
the identity.

\paragraph{\Cref{ex:tr-associativity-free}.}
By \cref{thm:tr-main}, \(\boxplus = \oplus\), and \(\oplus\) is associative because
\(\Phi\) carries it to addition on \(\C/\Lambda\), which is associative since
addition in \(\C\) is: \((P \oplus Q) \oplus S = \Phi((u+v)+w) = \Phi(u+(v+w)) =
P \oplus (Q \oplus S)\). It reduces precisely to associativity of \(+\) in \(\C\).
The direct geometric proof is hard because associativity of the chord--tangent rule
is a nonobvious incidence statement about nine points and lines (the
Cayley--Bacharach/\(9\)-point theorem); the torus makes it a triviality by
linearizing the group.

\paragraph{\Cref{ex:tr-intrinsic}.}
By \cref{prop:tr-intrinsic} the operation \(\boxplus\) is defined using only lines
and \(O\) on \(E\), so it is one and the same regardless of how \(E\) was
presented. \Cref{thm:tr-main} applied to \(\Phi\) gives \(\oplus = \boxplus\), and
applied to \(\Phi'\) gives \(\oplus' = \boxplus\); hence \(\oplus = \oplus'\). Since
both \(\Phi\) and \(\Phi'\) are isomorphisms onto the same \((E, \boxplus)\), the
composite \(\Phi^{-1}\!\circ \Phi'\) is a group isomorphism
\(\C/\Lambda' \to \C/\Lambda\), forcing \(\Lambda\) and \(\Lambda'\) to be
homothetic (\cref{prop:lat-normal-form}): the same cubic comes only from
homothetic lattices.

\paragraph{\Cref{ex:tr-translation}.}
Under \(\Phi\), the translation \(P \mapsto P \oplus A\) corresponds to
\(z \mapsto z + a\) on \(T\), where \(a = \Phi^{-1}(A)\); this is the biholomorphism
\(\tau_a\) of \cref{prop:tg-lie}. Conjugating by \(\Phi\),
\(P \mapsto P \oplus A = \Phi(\Phi^{-1}(P) + a)\) is a composition of
biholomorphisms, hence a biholomorphism of \(E\). The addition map \(E \times E \to
E\) and inversion are holomorphic (transported from \(T\)), so \(E\) is a complex
Lie group acting on itself by holomorphic translations.

\section{Chapter 12}
\label{sec:soln-ch12}

\paragraph{\Cref{ex:cases-chord-vieta}.}
Substituting \(y = mx + k\) into \(y^2 = 4x^3 - g_2 x - g_3\) gives
\(4x^3 - m^2 x^2 - (g_2 + 2mk)x - (g_3 + k^2) = 0\), a cubic whose three roots are
\(x_P, x_Q, x_R\). By Vieta's formula the sum of the roots is the negative of the
\(x^2\)-coefficient over the leading coefficient: \(x_P + x_Q + x_R = m^2/4\),
hence \(x_R = \tfrac{m^2}{4} - x_P - x_Q\). The third intersection is
\(R = (x_R, m x_R + k)\), and \(P \oplus Q = -R = (x_R, -(m x_R + k))\).

\paragraph{\Cref{ex:cases-tangent-slope}.}
Differentiating \(y^2 = 4x^3 - g_2 x - g_3\) implicitly: \(2y\,y' = 12x^2 - g_2\),
so the tangent slope at \(P = (x_P, y_P)\) is \(m = (12 x_P^2 - g_2)/(2 y_P)\). The
tangent meets \(E\) at \(P\) with multiplicity two, so in the cubic of the previous
exercise \(x_P\) is a double root and \(x_P + x_P + x_{2P} = m^2/4\), giving
\(x_{2P} = \tfrac{m^2}{4} - 2x_P\). If \(y_P = 0\), the slope \(m\) is infinite:
the tangent is vertical, meeting \(E\) at \(P\) (doubly) and \(O\), so
\(R = O\) and \(2P = \ominus O = O\), matching that \(P\) is two-torsion.

\paragraph{\Cref{ex:cases-two-torsion-sum}.}
The line \(y = 0\) (the \(x\)-axis) meets \(E : y^2 = 4x^3 - g_2 x - g_3\) where
\(4x^3 - g_2 x - g_3 = 0\), i.e.\ at \((e_1,0), (e_2,0), (e_3,0) = T_1, T_2, T_3\);
these are exactly the nonzero two-torsion points. As three points on a line,
\cref{thm:lin-collinear} gives \(T_1 \oplus T_2 \oplus T_3 = O\). In
\(E[2] \cong (\Z/2\Z)^2\) the three nonzero elements are \(T_1, T_2, T_3\), and
indeed the sum of all nonzero elements of \((\Z/2\Z)^2\) is \((1,1)+(1,0)+(0,1) =
(0,0)\), consistent.

\paragraph{\Cref{ex:cases-flex-three-torsion}.}
If \(3P = O\), then in the divisor language \(3[P] - 3[O]\) is principal, so by
\cref{thm:lin-collinear} (with \(P = Q = R\)) the line \(L\) with \(L \cdot E = 3P\)
exists --- the tangent at \(P\) meets \(E\) only at \(P\), with multiplicity three,
which is the definition of an inflection point. Conversely a flex \(P\) has
\(L \cdot E = 3P\), so \(3P \oplus = O\), i.e.\ \(3P = O\). The point \(O\) is a
flex by \cref{prop:emb-flex} (the line at infinity meets \(E\) at \(O\) with
multiplicity three), consistent with \(3O = O\).

\section{Chapter 13}
\label{sec:soln-ch13}

\paragraph{\Cref{ex:coord-yform}.}
With \(b = y_1 - m x_1\), \(-(m x_3 + b) = -m x_3 - (y_1 - m x_1) = m(x_1 - x_3)
- y_1\). The second form is preferable numerically because it uses the input point
\((x_1, y_1)\) directly and avoids forming the intercept \(b\) as a separate
quantity (one fewer subtraction, less cancellation error); it also makes the
geometric reading clear --- start at \(P\), move by slope \(m\) over the run
\(x_1 - x_3\), then the sign records the reflection.

\paragraph{\Cref{ex:coord-commute}.}
The slope \(m = (y_2-y_1)/(x_2-x_1) = (y_1-y_2)/(x_1-x_2)\) is unchanged by
swapping \(P \leftrightarrow Q\). Then \(x_3 = m^2/4 - x_1 - x_2\) is symmetric in
\(x_1, x_2\). For \(y_3\), note the line intercept \(b = y_1 - m x_1 = y_2 - m x_2\)
is the same computed from either point (both lie on the line), so
\(y_3 = -(m x_3 + b)\) is manifestly symmetric. Hence \(P \oplus Q = Q \oplus P\).

\paragraph{\Cref{ex:coord-limit}.}
On \(E\), \(y_2^2 - y_1^2 = 4(x_2^3 - x_1^3) - g_2(x_2 - x_1)\). Factoring,
\((y_2-y_1)(y_2+y_1) = (x_2-x_1)\bigl(4(x_2^2 + x_2 x_1 + x_1^2) - g_2\bigr)\), so
for \(x_2 \ne x_1\),
\[
  m = \frac{y_2 - y_1}{x_2 - x_1} = \frac{4(x_2^2 + x_2 x_1 + x_1^2) - g_2}{y_2 + y_1}.
\]
As \(Q \to P\) (so \(x_2 \to x_1\), \(y_2 \to y_1 \ne 0\)) this tends to
\((4 \cdot 3x_1^2 - g_2)/(2y_1) = (12x_1^2 - g_2)/(2y_1)\), the tangent slope.
Correspondingly Vieta's \(x_1 + x_2 + x_R = m^2/4\) becomes \(2x_1 + x_R = m^2/4\),
so the addition formula limits to the doubling formula.

\paragraph{\Cref{ex:coord-vertical}.}
If \(x_1 = x_2\) then \(y_1^2 = y_2^2\) (both equal \(4x_1^3 - g_2 x_1 - g_3\)), so
\(y_2 = \pm y_1\); since \(P \ne Q\), \(y_2 = -y_1\), i.e.\ \(Q = \ominus P\). The
slope \((y_2-y_1)/(x_2-x_1)\) has zero denominator and is undefined --- the line is
vertical --- and the third intersection is \(O\), so \(P \oplus Q = \ominus O = O\).
Pseudocode:
\begin{quote}\ttfamily
add(P, Q):\\
\ \ if P = O: return Q;\quad if Q = O: return P\\
\ \ if x1 = x2:\\
\ \ \ \ if y1 = -y2: return O\quad // includes y1 = y2 = 0\\
\ \ \ \ else: return double(P)\quad // here P = Q, y1 != 0\\
\ \ else:\\
\ \ \ \ m = (y2-y1)/(x2-x1)\\
\ \ \ \ x3 = m\^{}2/4 - x1 - x2\\
\ \ \ \ return (x3, m*(x1-x3)-y1)
\end{quote}

\paragraph{\Cref{ex:coord-check}.}
With \(x_{2P} = \tfrac{25}{24}\) and \(y_{2P} = -\tfrac{35}{24\sqrt6}\):
\(y_{2P}^2 = \tfrac{1225}{576 \cdot 6} = \tfrac{1225}{3456}\), while
\(4x_{2P}^3 - 4x_{2P} = 4\cdot\tfrac{15625}{13824} - \tfrac{100}{24}
= \tfrac{62500}{13824} - \tfrac{57600}{13824} = \tfrac{4900}{13824} = \tfrac{1225}{3456}\).
They agree, so \(2P \in E\). Its inverse is the reflection
\(\ominus(2P) = \bigl(\tfrac{25}{24}, +\tfrac{35}{24\sqrt6}\bigr)\).

\section{Chapter 14}
\label{sec:soln-ch14}

\paragraph{\Cref{ex:aa-h-order}.}
\(h = \wp' - A\wp - B\) is a \(\C\)-linear combination of \(\wp', \wp, 1\), hence
elliptic. Its poles can only be at lattice points; at \(0\), \(\wp'\) has a pole
of order three and \(\wp\) of order two, so \(h\) has a pole of order exactly three
(the \(\wp'\) term dominates). Thus \(h\) has order three, and by
\cref{thm:ell-valence} exactly three zeros. By the choice of \(A, B\), \(h(u) =
\wp'(u) - A\wp(u) - B = 0\) and likewise \(h(v) = 0\), so \(u, v\) are zeros. The
case \(u \equiv \pm v\) is excluded because then \(\wp(u) = \wp(v)\) and the slope
\(A\) has a zero denominator (for \(u \equiv v\) one must instead take the limit,
giving the duplication formula; for \(u \equiv -v\) the sum \(u+v \equiv 0\) is a
pole of \(\wp\)).

\paragraph{\Cref{ex:aa-determinant-expand}.}
Expanding along the first column, the determinant equals
\[
  \bigl[\wp(v)(-\wp'(u{+}v)) - \wp(u{+}v)\wp'(v)\bigr]
  - \bigl[\wp(u)(-\wp'(u{+}v)) - \wp(u{+}v)\wp'(u)\bigr]
  + \bigl[\wp(u)\wp'(v) - \wp(v)\wp'(u)\bigr] = 0,
\]
a linear equation in the row \((1, \wp(u+v), -\wp'(u+v))\). Geometrically it says
\((\wp(u+v), -\wp'(u+v))\) lies on the line through \(P, Q\); its \(x\)-coordinate
is therefore the third root \(\tfrac{A^2}{4} - \wp(u) - \wp(v)\) of the
intersection cubic, which is \(\wp(u+v)\) --- exactly \cref{thm:aa-addition}.

\paragraph{\Cref{ex:aa-duplication-match}.}
Differentiating \(\wp'^2 = 4\wp^3 - g_2\wp - g_3\) gives \(2\wp'\wp'' = (12\wp^2 -
g_2)\wp'\), so \(\wp'' = 6\wp^2 - \tfrac12 g_2\). Hence
\[
  \frac{\wp''(u)}{\wp'(u)} = \frac{6\wp(u)^2 - \tfrac12 g_2}{\wp'(u)}
  = \frac{12\wp(u)^2 - g_2}{2\wp'(u)} = \frac{12x^2 - g_2}{2y},
\]
with \(x = \wp(u)\), \(y = \wp'(u)\) --- precisely the tangent slope of
\cref{thm:coord-double}. So the analytic duplication slope and the geometric
tangent slope agree.

\paragraph{\Cref{ex:aa-symmetry}.}
Swapping \(u \leftrightarrow v\) negates both \(\wp'(u) - \wp'(v)\) and
\(\wp(u) - \wp(v)\), leaving the slope \(A\) and hence \(A^2\) unchanged; the term
\(-\wp(u) - \wp(v)\) is symmetric, so \cref{thm:aa-addition} is symmetric. As
\(v \to 0\), \(\wp(v) \sim v^{-2} \to \infty\) and \(\wp'(v) \sim -2v^{-3}\), so
\(A \sim \wp'(v)/(-\wp(v)) \sim (-2v^{-3})/(-v^{-2}) = 2v^{-1}\) and
\(A^2/4 \sim v^{-2} \sim \wp(v)\); the divergent \(-\wp(v)\) and
\(+A^2/4\) cancel to leading order, leaving \(\wp(u+0) = \wp(u)\).

\paragraph{\Cref{ex:aa-three-term}.}
Replacing \(v\) by \(-v\) in \cref{thm:aa-addition} and using \(\wp(-v) = \wp(v)\),
\(\wp'(-v) = -\wp'(v)\),
\[
  \wp(u-v) = -\wp(u) - \wp(v) + \tfrac14\Bigl(\tfrac{\wp'(u) + \wp'(v)}{\wp(u) - \wp(v)}\Bigr)^2 .
\]
Adding to the formula for \(\wp(u+v)\), the two square terms combine via
\((a-b)^2 + (a+b)^2 = 2a^2 + 2b^2\) with \(a = \wp'(u)\), \(b = \wp'(v)\):
\[
  \wp(u+v) + \wp(u-v) = -2\wp(u) - 2\wp(v)
  + \frac{\wp'(u)^2 + \wp'(v)^2}{2(\wp(u) - \wp(v))^2}.
\]
By the differential equation \(\wp'(u)^2 = 4\wp(u)^3 - g_2\wp(u) - g_3\) (and
likewise for \(v\)), the numerator is a polynomial in \(\wp(u), \wp(v)\), so the
right-hand side is rational in \(\wp(u), \wp(v)\).

\section{Chapter 15}
\label{sec:soln-ch15}

\paragraph{\Cref{ex:wk-numeric-check}.}
\(\bigl(\tfrac{28\sqrt6}{125}\bigr)^2 = \tfrac{784 \cdot 6}{15625} =
\tfrac{4704}{15625}\). On the other side, with \(x_3 = -\tfrac{24}{25}\),
\(4x_3^3 - 4x_3 = 4\bigl(-\tfrac{13824}{15625}\bigr) + \tfrac{96}{25}
= -\tfrac{55296}{15625} + \tfrac{60000}{15625} = \tfrac{4704}{15625}\). They agree,
so \(P \oplus Q \in E\). Its inverse is \(\ominus(P\oplus Q) =
\bigl(-\tfrac{24}{25}, +\tfrac{28\sqrt6}{125}\bigr)\).

\paragraph{\Cref{ex:wk-doubling-symbolic}.}
The tangent at \(P\) has slope \(m = (12x_1^2 - g_2)/(2y_1)\) and intercept
\(b = y_1 - m x_1\). It meets \(E\) at \(P\) doubly, so \(x_1\) is a double root of
\(4x^3 - m^2 x^2 - (g_2 + 2mb)x - (g_3 + b^2) = 0\); matching against
\(4(x - x_1)^2(x - x_3)\), the \(x^2\)-coefficient gives \(2x_1 + x_3 = m^2/4\),
so \(x_3 = \tfrac{m^2}{4} - 2x_1\). Reflecting the third point,
\(2P = (x_3, -(m x_3 + b)) = \bigl(\tfrac{m^2}{4} - 2x_1,\, -m(\tfrac{m^2}{4} -
2x_1) - b\bigr)\), as in \cref{thm:coord-double}.

\paragraph{\Cref{ex:wk-standard-recompute}.}
By \cref{prop:wk-standard} with \(g_2 = 4\), \(g_3 = 0\): \(a = -1\), \(b = 0\), so
\(E' : Y^2 = X^3 - X\). In these coordinates \(P = (2, \sqrt6)\),
\(Q = (-\tfrac12, \tfrac{\sqrt6}{4})\) (halving the \(y\)'s), and the slope is
\(M = m/2 = \tfrac{3\sqrt6}{10}\), \(M^2 = \tfrac{54}{100}\). Then
\(X_3 = M^2 - X_1 - X_2 = \tfrac{27}{50} - \tfrac32 = -\tfrac{24}{25}\) (the same
\(x\)) and \(Y_3 = M(X_1 - X_3) - Y_1 = -\tfrac{14\sqrt6}{125}\). Mapping back by
\(y = 2Y\) gives \(y_3 = -\tfrac{28\sqrt6}{125}\), matching \cref{sec:wk-numeric}.

\paragraph{\Cref{ex:wk-pitfall-instances}.}
\emph{Pitfall 1.} Compute \(P \oplus (\ominus P)\) for \(P = (2, 2\sqrt6)\): the
line is the vertical \(x = 2\), which has \emph{no} third finite intersection. If
\(O\) is forgotten the sum looks undefined; including \(O\), the line meets \(E\)
at \(P, \ominus P, O\), so \(P \oplus (\ominus P) = O\).
\emph{Pitfall 2.} In \cref{sec:wk-numeric} the third intersection is
\(R = (-\tfrac{24}{25}, +\tfrac{28\sqrt6}{125})\); reporting \(R\) as \(P \oplus Q\)
gives the reflection of the true answer --- wrong by a sign in \(y\). The correct
sum is \(\ominus R = (-\tfrac{24}{25}, -\tfrac{28\sqrt6}{125})\).
\emph{Pitfall 7.} On the singular cubic \(y^2 = 4x^3\) (\(g_2 = g_3 = 0\),
\(\Delta = 0\)) the origin \((0,0)\) is a cusp; the tangent there is not defined and
the doubling formula divides by \(y = 0\). With \(\Delta \ne 0\)
(\cref{cor:wp-smooth}) no such point exists and every step is well defined.

\paragraph{\Cref{ex:wk-associativity-numeric}.}
Take \(S = (-1, 0)\), a two-torsion point. From \cref{sec:wk-numeric},
\(P \oplus Q = (-\tfrac{24}{25}, -\tfrac{28\sqrt6}{125})\). Computing with the
formulas (slope cases handled as in \cref{ex:coord-vertical}):
\[
  (P \oplus Q) \oplus S = (49, 280\sqrt6) = P \oplus (Q \oplus S),
\]
so the two bracketings agree, confirming associativity numerically (guaranteed by
\cref{rem:tr-associativity}). Both intermediate sums lie on \(E\), as a direct
substitution checks.

\section{Chapter 16}
\label{sec:soln-ch16}

\paragraph{\Cref{ex:syn-j1728}.}
With \(g_2 = 4\), \(g_3 = 0\): \(\Delta = g_2^3 - 27 g_3^2 = 64 - 0 = 64\), and
\[
  j = 1728\,\frac{g_2^3}{\Delta} = 1728\,\frac{64}{64} = 1728 .
\]
For a curve with \(g_2 = 0\) (and \(g_3 \ne 0\)), \(j = 1728 \cdot 0/\Delta = 0\):
these are the curves with \(j = 0\), having complex multiplication by
\(\Z[\zeta_3]\).

\paragraph{\Cref{ex:syn-torsion}.}
A point \([z]\) satisfies \(n[z] = [0]\) iff \(nz \in \Lambda\) iff
\(z \in \tfrac1n\Lambda\). Hence \(E[n] = \tfrac1n\Lambda/\Lambda\). The map
\(\tfrac1n\Lambda/\Lambda \to \Lambda/n\Lambda\), \(z \mapsto nz\), is an
isomorphism, and \(\Lambda \cong \Z^2\) gives \(\Lambda/n\Lambda \cong (\Z/n\Z)^2\).
The map \([n]: [z] \mapsto [nz]\) has fibers the cosets of \(E[n]\), of which there
are \(|E[n]| = n^2\); so \([n]\) is \(n^2\)-to-one and has degree \(n^2\).

\paragraph{\Cref{ex:syn-weight}.}
From \(G_k(\alpha\Lambda) = \sum_{\lambda\ne0}(\alpha\lambda)^{-k} =
\alpha^{-k} G_k(\Lambda)\): \(g_2 = 60 G_4 \mapsto \alpha^{-4} g_2\) and
\(g_3 = 140 G_6 \mapsto \alpha^{-6} g_3\). Then \(g_2^3 \mapsto \alpha^{-12} g_2^3\),
\(g_3^2 \mapsto \alpha^{-12} g_3^2\), so \(\Delta = g_2^3 - 27 g_3^2 \mapsto
\alpha^{-12}\Delta\). Finally \(j = 1728 g_2^3/\Delta \mapsto
1728\,\alpha^{-12}g_2^3 / (\alpha^{-12}\Delta) = j\): invariant. In modular-forms
language \(g_2, g_3, \Delta, j\) have weights \(4, 6, 12, 0\).

\paragraph{\Cref{ex:syn-hasse}.}
For \(E : y^2 = x^3 + 1\) over \(\F_5\), the squares are \(\{0,1,4\}\). Checking
each \(x\): \(x=0 \Rightarrow y^2 = 1\), \(y = \pm1\); \(x=1 \Rightarrow y^2 = 2\),
none; \(x=2 \Rightarrow y^2 = 9 = 4\), \(y = \pm2\); \(x=3 \Rightarrow y^2 = 28 =
3\), none; \(x=4 \Rightarrow y^2 = 65 = 0\), \(y = 0\). The affine points are
\((0,1), (0,4), (2,2), (2,3), (4,0)\) --- five of them --- plus \(O\), so
\(\#E(\F_5) = 6\). Then \(a_5 = 5 + 1 - 6 = 0\) and \(|a_5| = 0 \le 2\sqrt5 \approx
4.47\), satisfying Hasse. The point \((4,0)\) has order \(2\); a group of order \(6\)
with an element of order \(2\) is \(\Z/6\Z\), so \(E(\F_5) \cong \Z/6\Z\).
